
\documentclass{article}
\usepackage{colortbl}
\usepackage{makecell}
\usepackage{multirow}
\usepackage{supertabular}

\begin{document}

\newcounter{utterance}

\centering \large Interaction Transcript for game `referencegame', experiment `line\_grids\_columns', episode 11 with reka{-}flash{-}3.1{-}t0.0.
\vspace{24pt}

{ \footnotesize  \setcounter{utterance}{1}
\setlength{\tabcolsep}{0pt}
\begin{supertabular}{c@{$\;$}|p{.15\linewidth}@{}p{.15\linewidth}p{.15\linewidth}p{.15\linewidth}p{.15\linewidth}p{.15\linewidth}}
    \# & $\;$A & \multicolumn{4}{c}{Game Master} & $\;\:$B\\
    \hline

    \theutterance \stepcounter{utterance}  
    & & \multicolumn{4}{p{0.6\linewidth}}{
        \cellcolor[rgb]{0.9,0.9,0.9}{
            \makecell[{{p{\linewidth}}}]{
                \texttt{\tiny{[P1$\langle$GM]}}
                \texttt{You are given three grids, where each of them is 5 by 5 in size.} \\
\texttt{Grids have empty cells marked with "▢" and filled cells marked with "X".} \\
\texttt{Your task is to generate a referring expression that best describes the target grid while distinguishing it from the two other distractor grids.} \\
\texttt{The first grid is the target grid, and the following two grids are the distractors.} \\
\\ 
\texttt{Target grid:} \\
\\ 
\texttt{X X X □ □} \\
\texttt{X X X □ □} \\
\texttt{X X X □ □} \\
\texttt{X X X □ □} \\
\texttt{X X X □ □} \\
\\ 
\texttt{Distractor grid 1:} \\
\\ 
\texttt{□ X X X □} \\
\texttt{□ X X X □} \\
\texttt{□ X X X □} \\
\texttt{□ X X X □} \\
\texttt{□ X X X □} \\
\\ 
\texttt{Distractor grid 2:} \\
\\ 
\texttt{□ X X □ X} \\
\texttt{□ X X □ X} \\
\texttt{□ X X □ X} \\
\texttt{□ X X □ X} \\
\texttt{□ X X □ X} \\
\\ 
\texttt{Instruction: Describe the target grid.} \\
\texttt{Generate the referring expression starting with the tag "Expression: " for the given target grid. Omit any other text.} \\
            }
        }
    }
    & \\ \\

    \theutterance \stepcounter{utterance}  
    & \multicolumn{4}{p{0.6\linewidth}}{
        \cellcolor[rgb]{0.9,0.9,0.9}{
            \makecell[{{p{\linewidth}}}]{
                \texttt{\tiny{[P1$\rangle$GM]}}
                \texttt{<reasoning>} \\
\texttt{Let's tackle this problem. I need to create a referring expression that describes the target grid and distinguishes it from the two distractor grids. I'll start by carefully looking at each grid to understand their structures.} \\
\\ 
\texttt{First, the target grid. All rows have three X's followed by two □'s. So each row is "X X X □ □". That's five columns. So the pattern is three filled cells in a row, then two empty. And this is consistent across all five rows. So the entire grid is a 3x2 block of X's? Wait, no. Wait, each row has three X's, but they are in the same positions each row. So rows 1{-}5 each have X's in columns 1{-}3, and □ in 4{-}5. So the X's form a 5x3 rectangle on the left side. So the target grid is a solid block of X's occupying the first three columns of all five rows.} \\
\\ 
\texttt{Now, Distractor grid 1: Each row starts with a □, then three X's, then a □. So each row is "□ X X X □". So the X's are in columns 2{-}4 of each row. So this is a horizontal block of three X's in the middle of each row, but shifted to the right compared to the target. So the X's here are in columns 2{-}4, whereas the target is columns 1{-}3. So the position of the X's is different.} \\
\\ 
\texttt{Distractor grid 2: Each row is "□ X X □ X". So here, the X's are in columns 2 and 3, and 5. Wait, let me check. The first row is □ X X □ X. So columns 2 and 3 are X's, then column 5 is X. So each row has two X's in the middle and one at the end. So the pattern here is two X's in columns 2{-}3 and one in column 5. So it's not a solid block. So the distractor 2 has a different arrangement.} \\
\\ 
\texttt{So the target is a solid 3x5 block on the left. Distractor 1 is a solid 3x5 block shifted right by one column. Distractor 2 is a different pattern with X} \\
            }
        }
    }
    & & \\ \\

    \theutterance \stepcounter{utterance}  
    & & & \multicolumn{2}{p{0.3\linewidth}}{
        \cellcolor[rgb]{0.9,0.9,0.9}{
            \makecell[{{p{\linewidth}}}]{
                \texttt{\tiny{[GM$|$GM]}}
                \texttt{Invalid generated expression} \\
            }
        }
    }
    & & \\ \\

\end{supertabular}
}

\end{document}
